\documentclass[10pt,a4paper]{amsart}
\usepackage[margin=19mm]{geometry}
\usepackage{amsmath,amssymb,mathtools}
\usepackage[scr=rsfs,cal=euler]{mathalfa}
\usepackage{xfrac}
\usepackage[unicode,hidelinks,bookmarks=false]{hyperref}
\newcommand{\ie}{i.e.\ }
\newcommand{\cf}{cf.\ }
\newcommand{\resp}{respectively\ }
\newcommand{\Subsection}{\subsection}
\providecommand{\nobreakdash}{\nobreak\hbox{-}}
\setlength{\emergencystretch}{2em}
\multlinegap0pt
\usepackage{amssymb}
\usepackage{mathtools}
\usepackage[shortlabels]{enumitem}
\newtheorem{claim}{Claim}
\usepackage{cleveref}
\crefname{claim}{Claim}{Claims}
\newcommand{\M}{\mathcal M}
\newcommand{\Rcal}{\mathcal R}
\newcommand{\X}{\mathcal X}
\title{Extremal Families for the Erd\H{o}s--Kleitman Problem: The Missing Constructions}
\author{Cheng Chi, Yan Wang}
\date{Source: arXiv:2607.25611v2 (CC BY 4.0)}
\begin{document}
\maketitle

\noindent\emph{Source and licence.} Extremal Families for the Erd\H{o}s--Kleitman Problem: The Missing Constructions. Version arXiv:2607.25611v2 (CC BY 4.0), distributed by arXiv under the Creative Commons Attribution 4.0 International licence, http://creativecommons.org/licenses/by/4.0/, as recorded in the arXiv licence metadata for this exact version and shown on the abstract page https://arxiv.org/abs/2607.25611v2. Full licence notice: \texttt{licenses/arxiv-2607.25611-CC-BY-4.0.txt}.\par\smallskip

\noindent\textit{Editorial scope.} Counted unit: the article's claim cl:defect-switch-admissible-exists (source0.tex lines 1683--1686). The article's complete original proof of it is delivered with it (source0.tex lines 1688--1819, one \texttt{proof} environment of 132 lines). No mathematics is written, repaired or re-derived: the statement and the proof are the article's own lines, in the article's own notation, and no retained line is substituted or re-flowed.  Retained uncounted as the source-local context the proof needs: lines 1109--1111; lines 1161--1167; lines 1168--1170; lines 1670--1674.  Nothing was omitted from any retained span, so the package carries no omission record.  No retained line contains a citation, so the wrapper carries no bibliography and no bibliography entry.  No retained line contains a figure environment, an image-inclusion command or drawing code, so no image is delivered and the package declares no asset dependency.  Editorial formatting only, in addition: an AMS \texttt{amsart} preamble that restates the article's own declarations and macros used by the retained lines, namely the environments \texttt{claim} on separate counters, together with the standard packages those lines need.  The retained environments print in this document as Claim 1 in order of appearance, whereas the article prints the same items with the numbering of its own full document.  The complete native source of the article as distributed in the arXiv e-print for this exact version is bundled unchanged as \texttt{sources/206196/source0.tex}.
\begin{claim}\label{cl:defect-coarse-bounds}
        There is a constant $C_0=C_0(m,k)$ such that $|\X(A)|\le C_0cs^{m-1}$ and $d\le C_0c$.
\end{claim}

\begin{claim}\label{cl:defect-capture}
        We have $|\X(A)\setminus X_U^\star|\le M_1(\nu(\Rcal_U)+d)E_k(s,c)$.
        Moreover, if $U=\emptyset$, then
        \[
                |\X(A)|\le M_1\nu(\X(A))B_k(s,c).
        \]
\end{claim}

\begin{claim}\label{cl:defect-completion}
        If $\nu(\Rcal_U)>d$, then $|\M(A)\setminus Z_U|\ge M_3(\nu(\Rcal_U)-d)B_k(s,c)$.
\end{claim}

\begin{equation}\label{eq:defect-switch-admissible}
        |\X(A)\setminus T(S)|
        \le
        |\M(A)\setminus Z_U|+M_2dE_k(s,c).
\end{equation}

\begin{claim}\label{cl:defect-switch-admissible-exists}
        There exist $U'\subseteq U$ and
        $S\subseteq R$ such that $(U',S)$ is switch-admissible.
\end{claim}

\begin{proof}
        Let $\mu=\nu(X_U^\star)$, and choose a maximum matching $\mathcal J_0
                =
                \{P_r\cup\{z_r\}:1\le r\le\mu\}
                \subseteq X_U^\star$.
        Let $S_\mu=\{z_1,\ldots,z_\mu\}$.
        Since $X_U^\star\subseteq\Rcal_U$, we have
        $\mu\le\nu(\Rcal_U)$.
        We split the proof into two cases based on the value of $d$.

        \medskip
        \noindent\textbf{Case 1.} Suppose that $d>|R|$.
        \medskip

        In this case, choose $U'\subseteq U$ with $|U'|=|R|$, and take $S=R$.
        Thus, $X_U^\star\subseteq T(S)$, because every member of
        $X_U^\star$ has the form $P\cup\{z\}$ with
        $P\in\binom{A_0}{m-1}$ and $z\in R=S$.

        Every member of $\Rcal_U$ is $A$-extra, and hence meets $R$.
        Thus, any matching in $\Rcal_U$ uses distinct vertices of $R$,
        so $\nu(\Rcal_U)\le |R|<d$.
        By \cref{cl:defect-capture},
        \[
                |\X(A)\setminus T(S)|
                \le
                |\X(A)\setminus X_U^\star|
                \le
                M_1(\nu(\Rcal_U)+d)E_k(s,c)
                \le
                2M_1dE_k(s,c)\le M_2dE_k(s,c).
        \]
        The last inequality follows from the hierarchy.  Hence,
        \eqref{eq:defect-switch-admissible} holds.

        It remains to verify the size condition.
        Here $|U'|=|S|=|R|\ge1$.
        By \cref{cl:defect-coarse-bounds}, $d\le C_0c$.
        Since $|R|=bc+1\ge c$ and $M_2$ is chosen sufficiently large
        in the hierarchy, $d\le C_0c\le M_2|R|$.
        Thus, $(U',S)$ is switch-admissible in this case.

        \medskip
        \noindent\textbf{Case 2.} Suppose that $d\le |R|$.
        \medskip

        In this case, let $U'=U$.
        Choose $S\subseteq R$ with $|S|=d$ so that $|S\cap S_\mu|$ is
        maximal.
        Given a number $x$, write $x_+$ for $\max\{x,0\}$.
        Thus, $|S_\mu\setminus S|=(\mu-d)_+$ since $|S|=d$ and $|S_\mu|=\mu$ by definition.
        We first bound the size of $X_U^\star\setminus T(S)$.
        \begin{claim}\label{cl:clean-center-XU-loss}
                There is a constant $C_7=C_7(m,k)>0$ such that
                \[
                        |X_U^\star\setminus T(S)|
                        \le
                        (\mu-d)_+\binom{|A_0|}{m-1}
                        +C_7\mu cs^{m-2}.
                \]
        \end{claim}

        \begin{proof}
                Every member of $X_U^\star$ has the form $P\cup\{z\}$,
                where $P\in\binom{A_0}{m-1}$ and $z\in R$.
                If $z\in S$, then this set belongs to $T(S)$.
                Hence only vertices $z\in R\setminus S$ can contribute to
                $X_U^\star\setminus T(S)$.
                Since $S$ maximizes $|S\cap S_\mu|$ among all
                $d$-subsets of $R$, the number of vertices of $S_\mu$
                left outside $S$ is exactly $(\mu-d)_+$.
                The members with $z\in S_\mu\setminus S$ therefore
                contribute at most $(\mu-d)_+\binom{|A_0|}{m-1}$.

                It remains to count members
                $P\cup\{z\}\in X_U^\star\setminus T(S)$ with
                $z\notin S_\mu$.
                Let $Q=\bigcup_{r=1}^{\mu}P_r$.
                If such a member had $P\cap Q=\emptyset$, then
                $P\cup\{z\}$ would be disjoint from every member
                $P_r\cup\{z_r\}$ of $\mathcal J_0$: the sets $P$ and
                $P_r$ are disjoint by $P\cap Q=\emptyset$, the vertex
                $z$ is distinct from every $z_r$ because $z\notin S_\mu$,
                and $P\subseteq A_0$ while each $z_r\in R$.
                This would enlarge the matching $\mathcal J_0$,
                contradicting its maximality.
                Thus every remaining member has $P\cap Q\ne\emptyset$.

                We count these remaining members by first choosing
                $z\in R$, then a vertex of $P\cap Q$, and then the other
                $m-2$ vertices of $P$ from $A_0$.
                Since $|R|\le(b+1)c$, $|Q|\le(m-1)\mu$, and
                $|A_0|<(m+1)s$, their number is
                $O_{m,k}(\mu cs^{m-2})$.  Choosing $C_7=C_7(m,k)$
                sufficiently large proves the claim.
        \end{proof}

        We now combine this estimate with \cref{cl:defect-capture}.
        Since $\binom{|A_0|}{m-1}=O_{m,k}(s^{m-1})$,
        $\mu\le\nu(\Rcal_U)$, $cs^{m-2}\le E_k(s,c)$, and
        $B_k(s,c)=s^{m-1}+E_k(s,c)$, there is a constant
        $C_8=C_8(m,k,M_1)$ such that
        \[
                \begin{aligned}
                        |\X(A)\setminus T(S)|
                         & \le
                        |\X(A)\setminus X_U^\star|
                        +|X_U^\star\setminus T(S)| \\
                         & \le
                        C_8(\nu(\Rcal_U)-d)_+B_k(s,c)
                        +C_8dE_k(s,c)              \\
                         & \le
                        C_8(\nu(\Rcal_U)-d)_+B_k(s,c)
                        +M_2dE_k(s,c),
                \end{aligned}
        \]
        where the last inequality follows from the hierarchy.

        If $\nu(\Rcal_U)\le d$, then the first term vanishes, and
        \eqref{eq:defect-switch-admissible} follows immediately.
        If $\nu(\Rcal_U)>d$, then \cref{cl:defect-completion} gives
        $|\M(A)\setminus Z_U|\ge
                M_3(\nu(\Rcal_U)-d)B_k(s,c)$.
        Since $M_3$ is chosen after $M_1$ and $M_2$, the hierarchy
        ensures that $M_3>C_8$, and this again implies
        \eqref{eq:defect-switch-admissible}.

        Finally, $|U'|=|S|=d\ge1$ and
        $d\le M_2d=M_2|U'|$ are immediate in this case.
        Thus, $(U',S)$ is switch-admissible in Case 2.
        The two cases prove the claim.
\end{proof}

\end{document}
